% ModernLayout
\begingroup
\hypersetup{bookmarksdepth=3}
\tcbset{breakable/.default=false,after={\par\nobreak\smallskip}}
\newcommand{\ModernAnswerLabel}{{\bfseries\color{DeepBlue}解答}}
\RenewDocumentEnvironment{xesolution}{+b}{%
  \par\nobreak\smallskip\nobreak\begingroup\small
  \clubpenalties 2 10000 10000\displaywidowpenalties 2 10000 10000\noindent
  #1\par\endgroup\smallskip
}{}
\chapter{量子现象、原子与原子核物理}

\section{热辐射与能量量子}

黑体是在给定温度下对各波长辐射具有最大吸收率的理想系统。经典能量均分与连续电磁模式结合会在高频给出发散，
普朗克（Max Planck）假定频率为 $\nu$ 的模式与物质交换能量时以 $h\nu$ 为单位，得到正确的黑体谱。单位体积、单位频率间隔的谱能量密度为
\begin{equation}
 u_\nu(T)=\frac{8\pi h\nu^3}{c^3}\frac1{e^{h\nu/(k_BT)}-1}.
\end{equation}
低频极限 $h\nu\ll k_BT$ 时，令 $x=h\nu/(k_BT)$，利用 $e^x-1\simeq x$，得到
\[
 u_\nu\simeq\frac{8\pi\nu^2k_BT}{c^3}.
\]
这就是瑞利--金斯形式（Rayleigh--Jeans law）。高频端的指数因子压过三次幂，因而不存在经典理论的紫外发散。
对频率积分得到的是单位体积内的总辐射能量 $u=\int_0^\infty u_\nu\,\dd\nu$，单位为 $\mathrm{J/m^3}$，并不是功率。
各向同性辐射中，每单位立体角分配的能量密度为 $u/(4\pi)$。穿过单位面积的能流还要乘传播速度 $c$ 和法向投影 $\cos\theta$，因此出射半球给出
\[
 j=\int_0^{2\pi}\!\int_0^{\pi/2}\frac{cu}{4\pi}\cos\theta\sin\theta\,\dd\theta\,\dd\varphi
 =\frac{cu}{4}.
\]
同一角积分还说明单位面积向立体角元 $\dd\Omega$ 发射的功率满足 $\dd j/\dd\Omega=(j/\pi)\cos\theta$，即朗伯余弦定律；$\theta$ 从表面外法线量起。
所以黑体的总辐出度（单位面积的出射功率）为
\begin{equation}
 j=\sigma T^4,
\end{equation}
面积为 $S$ 的均匀黑体表面的出射功率为 $P=Sj$。这里 $j$ 的单位为 $\mathrm{W/m^2}$；若计算与温度为 $T_0$ 的黑体环境之间的净交换功率，则为 $\sigma S(T^4-T_0^4)$。
单位波长间隔的能量密度需作变量变换，而不能只把 $\nu$ 换成 $c/\lambda$：
\[
 u_\lambda=u_\nu(c/\lambda)\left|\frac{\dd\nu}{\dd\lambda}\right|
 =\frac{8\pi hc}{\lambda^5}\frac1{e^{hc/(\lambda k_BT)}-1},\qquad M_\lambda=\frac c4u_\lambda.
\]
令 $z=hc/(\lambda k_BT)$。固定温度时，$M_\lambda$ 的极大值对应 $z^5/(e^z-1)$ 的极大值，其导数为零要求
\[
 5(e^z-1)-ze^z=0,\qquad 5(1-e^{-z})=z.
\]
非零解 $z\simeq4.9651$ 给出维恩位移律（Wien's displacement law）$\lambda_{\max}T=b=hc/(4.9651k_B)$。
按频率和按波长表示的谱密度，其峰值不是简单的倒数关系。谱峰波长随温度升高向短波移动，但不能把
$hc/\lambda_{\max}$ 直接当作所有光子的平均能量。

热平衡时，同一频率和方向上的发射率等于吸收率，这就是基尔霍夫热辐射定律（Kirchhoff's law）。黑体的吸收率为1；实际物体的发射谱还必须乘上不超过1且依赖材料、波长和方向的发射率，不能只凭温度把所有表面都当成理想黑体。


\subsection{从谱积分到温度四次方}
在普朗克谱积分中令 $x=h\nu/(k_BT)$，则 $\dd\nu=(k_BT/h)\dd x$，所有温度因子可以提出积分号：
\[
 u=\frac{8\pi k_B^4T^4}{h^3c^3}
 \int_0^\infty\frac{x^3}{e^x-1}\,\dd x
 =\frac{8\pi^5k_B^4}{15h^3c^3}T^4.
\]
这里使用无量纲积分 $\int_0^\infty x^3/(e^x-1)\,\dd x=\pi^4/15$。结合 $j=cu/4$，得
$\sigma=2\pi^5k_B^4/(15h^3c^2)$。温度四次方既来自谱中的三次幂，也来自频率积分的一个尺度因子。

若球形天体半径为 $R$、距离为 $d\gg R$，忽略吸收且均匀辐射，则总光度 $L_*=4\pi R^2\sigma T^4$，探测器处辐照度
\[
 F=\frac{L_*}{4\pi d^2}=\sigma T^4\left(\frac Rd\right)^2.
\]
表面的辐出度与远处接收的辐照度同为 $\mathrm{W/m^2}$，却对应不同面积，不能因单位相同而直接相等。实际恒星大气、选择性吸收会改变谱形，本节只使用黑体近似。

\subsection{黑体谱的温标}

\begin{origexamplebox}{恒星表面温度与辐出度}
把一颗恒星表面近似为黑体，其单位波长间隔的辐射谱峰值在 $500\,\mathrm{nm}$。用
$b=2.898\times10^{-3}\,\mathrm{m\,K}$、$\sigma=5.67\times10^{-8}\,\mathrm{W\,m^{-2}K^{-4}}$，求表面温度和单位面积总辐射功率。
\end{origexamplebox}
\begin{xesolution}
\begin{center}\begin{minipage}{\linewidth}\centering\raggedright\ModernAnswerLabel\par\smallskip\centering
% Contract: common absolute 普朗克 scale; T=4000,6000 K; peaks b/T.
\begin{tikzpicture}
\begin{axis}[width=.83\linewidth,height=6cm,axis lines=left,
 xlabel={$\lambda\;(\mu\mathrm m)$},ylabel={$M_\lambda/(10^{13}\,\mathrm{W\,m^{-3}})$},
 xmin=.15,xmax=2.5,ymin=0,ymax=11.8,xtick={.5,1,1.5,2},ytick={0,2,4,6,8,10},
 legend style={draw=none,fill=none,at={(.97,.95)},anchor=north east},
 every axis label/.style={fill=none}]
 \addplot[DeepBlue,very thick,domain=.15:2.5,samples=300]{37.4177/(x^5*(exp(14387.77/(6000*x))-1))};
 \addlegendentry{$6000\,\mathrm K$}
 \addplot[TextGray,very thick,dashed,domain=.15:2.5,samples=300]{37.4177/(x^5*(exp(14387.77/(4000*x))-1))};
 \addlegendentry{$4000\,\mathrm K$}
 \draw[guide] (axis cs:.48296203365,0)--(axis cs:.48296203365,10.0053244776);
 \draw[guide] (axis cs:.72444305047,0)--(axis cs:.72444305047,1.31757359375);
\end{axis}
\end{tikzpicture}

\figcap{黑体单色辐出度随温度的变化：升温后总面积增大，峰值向短波移动。}\label{ch32:fig:blackbody}
\end{minipage}\end{center}


维恩位移律给出
\[
 T=\frac b{\lambda_{\max}}
 =\frac{2.898\times10^{-3}}{500\times10^{-9}}
 \approx\boxed{5.80\times10^3\,\mathrm K}.
\]
斯蒂芬--玻尔兹曼定律给出
\[
 j=\sigma T^4\approx5.67\times10^{-8}(5.80\times10^3)^4
 \approx\boxed{6.4\times10^7\,\mathrm{W/m^2}}.
\]
峰值波长乘温度的单位为 $\mathrm{m\,K}$；温度升高时 $\lambda_{\max}$ 变短且总功率按 $T^4$ 增大，与图~\ref{ch32:fig:blackbody} 一致。
\end{xesolution}



\section{光子、光电效应与康普顿散射}

光的量子具有
\begin{equation}
 E=h\nu=\frac{hc}{\lambda},\qquad p=\frac E c=\frac h\lambda.
\end{equation}
光电效应中，一个电子吸收一个光子的最简能量关系为
\begin{equation}
 K_{\max}=h\nu-W,
\end{equation}
其中 $W$ 为逸出功，只有 $h\nu\ge W$ 才能发生这一单光子光电过程。截止电压 $U_s$ 指使最快光电子也不能到达收集极的反向电压大小。
设发射极电势为零，收集极电势为 $-U_s$；电子电荷为 $-e$，到达收集极需要增加的电势能为 $(-e)(-U_s)=eU_s$。临界时末动能为零，所以 $eU_s=K_{\max}$。
频率决定单个光子的能量和截止电压，光强主要改变单位时间到达的光子数；在多光子或强场条件下还会出现更复杂过程。

\begin{origexamplebox}{光电效应截止电压}
金属逸出功 $W=2.2\,\mathrm{eV}$，用波长 $400\,\mathrm{nm}$ 的单色光照射。求光电子最大动能和截止电压。
\end{origexamplebox}
\begin{xesolution}
\ModernAnswerLabel\quad 使用 $hc\simeq1240\,\mathrm{eV\,nm}$：
\[
 E_\gamma=\frac{1240}{400}\,\mathrm{eV}=3.10\,\mathrm{eV}.
\]
所以
\[
 K_{\max}=3.10-2.20=0.90\,\mathrm{eV},
 \qquad \boxed{U_s=0.90\,\mathrm V}.
\]
这里第二个数值不是把能量直接改写成电压。电子伏特的定义是
$1\,\mathrm{eV}=e\times1\,\mathrm V$，因此
\[
 U_s=\frac{K_{\max}}e
 =\frac{0.90\times1.602\times10^{-19}\,\mathrm J}
 {1.602\times10^{-19}\,\mathrm C}=0.90\,\mathrm V.
\]
\begin{center}\begin{minipage}{\linewidth}\centering
% Contract: collector has lower potential; electron travels right and slows.
% Plates are electrodes, rays enter cathode from upper left; not a circuit schematic.
\begin{tikzpicture}[>=Latex,font=\small]
 \draw[very thick] (0,-.6)--(0,1.3) (5,-.6)--(5,1.3);
 \node[below=6pt,align=center] at (0,-.6){发射极K\\$\varphi_K=0$};
 \node[below=6pt,align=center] at (5,-.6){收集极A\\$\varphi_A=-U_s$};
 \draw[DeepBlue,thick,->] (-1.2,2)--(0,.7) node[pos=.22,left=3pt]{$h\nu$};
 \draw[DeepBlue,thick,->] (.45,.55)--(2.4,.55) node[midway,above=5pt]{$e^-$ 向右运动};
 \draw[->] (3,-.05)--(4.3,-.05) node[midway,below=4pt]{$\vect E$};
 \draw[->] (2.4,1.55)--(1.1,1.55) node[midway,above=4pt]{$\vect F=-e\vect E$};
\end{tikzpicture}

\figcap{截止电压的电势约定。K接较高电势，A接较低电势；图示电场和受力方向相反。}\label{ch32:fig:photoelectric-apparatus}
\end{minipage}\end{center}
如图~\ref{ch32:fig:photoelectric-apparatus}，电子向右时电场力向左，故减速。
从K到A电势降低 $U_s$，电子电势能却增加 $eU_s$；
最快电子恰好耗尽初动能时到达临界截止，收集极电势为 $-0.90\,\mathrm V$。
若增加同频光强，饱和光电流可增大，但理想模型中的截止电压不变。
\end{xesolution}

\subsection{光电效应的直线判据}

\begin{origexamplebox}{由截止电压数据求普朗克常数}
同一种金属在频率 $f_1=6.0\times10^{14}\,\mathrm{Hz}$ 时截止电压为 $0.50\,\mathrm V$，
在 $f_2=8.0\times10^{14}\,\mathrm{Hz}$ 时为 $1.33\,\mathrm V$。估算 $h$、阈频和逸出功。
\end{origexamplebox}
\begin{xesolution}
\begin{center}\begin{minipage}{\linewidth}\centering\raggedright\ModernAnswerLabel\par\smallskip\centering
\begin{tikzpicture}
\begin{axis}[width=.72\linewidth,height=4.7cm,axis lines=left,
  xlabel={$f\;(10^{14}\,\mathrm{Hz})$},ylabel={$U_s\;(\mathrm V)$},xmin=3.5,xmax=9,ymin=0,ymax=1.8,
  xtick={4.8,6,8},xticklabels={$f_0$,$f_1$,$f_2$},ytick={.5,1.33},
  every axis label/.style={fill=none}]
  \addplot[DeepBlue,very thick,domain=4.795180723:8.8]{.415*x-1.99};
  \addplot[only marks,mark=*,DeepBlue] coordinates{(6,.5)(8,1.33)};
  \draw[guide] (axis cs:6,0)--(axis cs:6,.5);
  \draw[guide] (axis cs:8,0)--(axis cs:8,1.33);
\end{axis}
\end{tikzpicture}
\figcap{截止电压对频率的直线关系 $U_s=(h/e)f-W/e$。斜率决定普朗克常数，横截距决定阈频。}\label{ch32:fig:photoelectric-line}
\end{minipage}\end{center}
图~\ref{ch32:fig:photoelectric-line} 标出的两测量点用于本题的斜率计算，不能用连线在纵轴上的截距误认逸出功本身。


由 $eU_s=hf-W$，直线斜率为
（把两个频率对应的等式相减，逸出功 $W$ 消去）
\[
 \frac he=\frac{1.33-0.50}{(8.0-6.0)\times10^{14}}
 =4.15\times10^{-15}\,\mathrm{V\,s}.
\]
取 $e=1.602\times10^{-19}\,\mathrm C$，
\[
 \boxed{h\approx6.65\times10^{-34}\,\mathrm{J\,s}}.
\]
阈频
\[
 f_0=f_1-\frac{U_{s1}}{h/e}
 \approx\boxed{4.80\times10^{14}\,\mathrm{Hz}}.
\]
以电子伏特表示，$W/e=(h/e)f_0\approx1.99\,\mathrm V$，故
$\boxed{W\approx1.99\,\mathrm{eV}}$。两组数据的斜率单位为 $\mathrm{V\,s}$，乘电荷后为 $\mathrm{J\,s}$。
\end{xesolution}

\subsection{散射中的能量与动量守恒}
康普顿散射（Compton scattering）中，取初始自由电子静止，入射光子沿 $x$ 轴，散射光子与它成角 $\theta$。
设两光子能量分别为 $E,E'$，末态电子总能量为 $E_e$。能量和两个方向的动量守恒分别给出
\[
 E+m_ec^2=E'+E_e,\qquad p_{ex}=\frac{E-E'\cos\theta}{c},\qquad
 p_{ey}=-\frac{E'\sin\theta}{c}.
\]
因此 $p_e^2c^2=E^2+E'^2-2EE'\cos\theta$。再用 $E_e^2=m_e^2c^4+p_e^2c^2$，代入 $E_e=m_ec^2+E-E'$，消去两边共同项：
\[
 2m_ec^2(E-E')=2EE'(1-\cos\theta),\qquad
 \frac1{E'}-\frac1E=\frac{1-\cos\theta}{m_ec^2}.
\]
最后代入 $E=hc/\lambda$，得到波长改变
\begin{equation}
 \lambda'-\lambda=\frac{h}{m_ec}(1-\cos\theta).
\end{equation}
它说明高能电磁辐射交换能量和动量时具有不可忽略的量子性。

\subsection{康普顿散射几何}

\begin{origexamplebox}{九十度康普顿散射}
波长 $\lambda=0.0710\,\mathrm{nm}$ 的X射线被近似静止的自由电子散射，光子偏转 $90^\circ$。
取电子康普顿波长 $h/(m_ec)=0.002426\,\mathrm{nm}$，求散射光波长及散射前后光子能量。
\end{origexamplebox}
\begin{xesolution}
\begin{center}\begin{minipage}{\linewidth}\centering\raggedright\ModernAnswerLabel\par\smallskip\centering
% Contract: 90-degree photon scattering; recoil direction from p_e=p-p'.
\begin{tikzpicture}[font=\small,>=Latex]
 \coordinate (O) at (0,0);
 \draw[vector] (-3.6,0)--(O) node[midway,above=6pt] {入射光子 $\lambda$};
 \draw[guide] (O)--(3.1,0);
 \draw[vector] (O)--(0,2.7) node[above] {散射光子 $\lambda'$};
 \draw[vector] (O)--(2.7,-2.6108) node[right=5pt] {$\vect p_e$};
 \fill (O) circle (1.5pt);
 \node[below left=8pt] at (O) {$O$};
 \draw[DeepBlue,->] (.7,0) arc[start angle=0,end angle=90,radius=.7];
 \node at (1.2,1.05) {$\theta=90^\circ$};
 \node[align=center] at (-2,-1.1) {电子初始静止\\散射点为 $O$};
\end{tikzpicture}

\figcap{康普顿散射中的动量矢量方向，箭长不表示动量大小。散射角 $\theta$ 是入射与散射光子方向的夹角。}\label{ch32:fig:compton}
\end{minipage}\end{center}
图~\ref{ch32:fig:compton} 对应本题的 $90^\circ$ 散射：散射光子向上带走动量，反冲电子便必须具有向下的动量分量，同时保有向右的动量分量。电子箭头的方向由 $\vect p_e=\vect p-\vect p'$ 决定，而不是任意画出。


$\theta=90^\circ$ 时 $1-\cos\theta=1$，故
\[
 \boxed{\lambda'=0.0710+0.002426=0.07343\,\mathrm{nm}}.
\]
用 $hc=1.23984\,\mathrm{keV\,nm}$：
\[
 E=\frac{1.23984}{0.0710}\approx17.46\,\mathrm{keV},
 \qquad
 E'=\frac{1.23984}{0.07343}\approx16.89\,\mathrm{keV}.
\]
保留未舍入数值作差，光子损失约 $0.577\,\mathrm{keV}$，转化为电子动能。$\theta=0$ 时波长不变；$\theta=\pi$ 时改变量达到两倍康普顿波长。
图~\ref{ch32:fig:compton} 的出射光子向上，故按前面的动量守恒式有
\[
 p_{ex}=E/c>0,\qquad p_{ey}=-E'/c<0,
 \qquad \tan\phi=\frac{|p_{ey}|}{p_{ex}}=\frac{E'}E\approx0.967.
\]
电子在入射方向下方约 $44.0^\circ$ 反冲，与图中的箭头一致。
还可作独立能量检查：$p_e^2c^2=E^2+E'^2$，所以
\[
 K_e=\sqrt{(m_ec^2)^2+E^2+E'^2}-m_ec^2\approx0.577\,\mathrm{keV}.
\]
它与 $E-E'$ 一致；若把两光子的能量差全部当作电子的总能量，就会漏掉电子原有的静止能量。
\end{xesolution}



\section{物质波与量子态}

德布罗意关系（de Broglie relation）把粒子动量与波长联系起来：
\begin{equation}
 \lambda=\frac hp.
\end{equation}
电子、中子、原子乃至更大复合粒子都能在适当相干条件下产生衍射。这里“波”指量子态的概率振幅具有相位并可叠加，
“粒子”指探测时能量、电荷等以局域事件交换。量子对象不是一颗始终沿唯一轨迹运动的经典小球，也不是在三维空间中可直接测量的经典物质波；
实验装置和测量问题决定哪些概率分布可被观察。

记约化普朗克常数为 $\hbar=h/(2\pi)$。位置与动量的不确定关系为
\begin{equation}
 \Delta x\,\Delta p_x\ge\frac\hbar2.
\end{equation}
这里 $\Delta x=\sqrt{\langle x^2\rangle-\langle x\rangle^2}$，$\Delta p_x$ 同样是标准差，而不是两次测量值之差。它是量子态允许的概率分布宽度关系，不等同于仪器粗糙造成的普通测量误差。把波包限制得更窄需要更宽的动量成分。

\begin{origexamplebox}{电子经电势差加速后的波长}
电子从静止经电势差 $U$ 加速，且 $eU\ll m_ec^2$。求德布罗意波长，并说明何时需相对论修正。
\end{origexamplebox}
\begin{xesolution}
\ModernAnswerLabel\quad 非相对论能量关系 $eU=p^2/(2m_e)$ 给出
\[
 \boxed{\lambda=\frac h{\sqrt{2m_eeU}}}.
\]
若 $eU$ 与 $m_ec^2=511\,\mathrm{keV}$ 不再相比很小，应使用
\[
 (pc)^2=(m_ec^2+eU)^2-m_e^2c^4
\]
求动量。展开右侧并取正根，得到
\[
 pc=\sqrt{2m_ec^2eU+(eU)^2},\qquad
 \lambda=\frac{h}{\sqrt{2m_eeU}\sqrt{1+eU/(2m_ec^2)}}.
\]
当 $eU/(m_ec^2)\to0$ 时第二个根式趋于1，回到前面的低压公式。低压近似的适用条件来自能量尺度比较，而不是只看电压数值是否“很大”。
\end{xesolution}

\begingroup\tcbset{breakable/.default=true}\begin{starread}{薛定谔方程与自旋}
非相对论单粒子量子态由波函数 $\psi(\vect r,t)$ 表示，满足
\[
 \ii\hbar\frac{\partial\psi}{\partial t}
 =\left[-\frac{\hbar^2}{2m}\nabla^2+V(\vect r,t)\right]\psi.
\]
$|\psi|^2$ 是位置概率密度，满足 $\int |\psi|^2\,\dd^3r=1$。定态问题化为哈密顿（Hamilton）算符的本征值问题；边界条件使束缚态能量离散。
自旋是粒子的内禀角动量，不等同于带电小球绕自身轴作经典旋转。电子自旋测量沿任意轴只得到
$\pm\hbar/2$，原子中的离散单电子能级与泡利不相容原理决定基本的壳层填充：每个空间轨道至多容纳两个自旋态不同的电子。自旋磁矩则参与磁场中的能级分裂等磁相互作用，不能把磁矩本身说成壳层填充的原因。
\end{starread}\endgroup
\subsection{由边界条件认识能级离散}
以一维无限深势阱 $0<x<a$ 为模型：阱内 $V=0$，阱外概率为零，波函数在两壁满足 $\psi(0)=\psi(a)=0$。定态方程化为
\[
 \psi''+k^2\psi=0,\qquad E=\frac{\hbar^2k^2}{2m}.
\]
第一边界排除余弦项，第二边界要求 $\sin(ka)=0$。非零波函数对应 $k_n=n\pi/a$，$n=1,2,\ldots$，归一化后
\[
 \psi_n(x)=\sqrt{\frac2a}\sin\frac{n\pi x}{a},\qquad
 E_n=\frac{n^2\pi^2\hbar^2}{2ma^2}=\frac{n^2h^2}{8ma^2}.
\]
$n=0$ 只给处处为零的函数，不能归一化，所以基态能量不为零。这个模型说明离散值来自波动方程与边界条件，而不是任意把经典能量挑几个数值留下。
对电子和 $a=1.00\,\mathrm{nm}$，$E_1\approx0.376\,\mathrm{eV}$，$E_2-E_1=3E_1\approx1.13\,\mathrm{eV}$。阱宽减半时所有能级增为四倍。
无限势垒是理想边界；有限势垒会出现穿透尾部，不能继续强令阱外波函数处处为零。

衍射计算则先由加速能量求 $p$ 和 $\lambda=h/p$，再使用晶面的相位条件。加速电压、波长、散射角属于同一计算链；晶面夹角与法线夹角互为余角，使用布拉格式前应明确测角方式。

\subsection{电子衍射}

\begin{origexamplebox}{低能电子的德布罗意波长与布拉格角}
电子从静止经 $54.0\,\mathrm V$ 加速后入射晶体，晶面间距为 $0.200\,\mathrm{nm}$。
电子波长可用下式计算（$\lambda$ 以纳米计，$U$ 以伏特计）：
\[
 \lambda(\mathrm{nm})\approx\frac{1.227}{\sqrt{U(\mathrm V)}}.
\]
求电子波长和一级布拉格角。
\end{origexamplebox}
\begin{xesolution}
\begin{center}\begin{minipage}{\linewidth}\centering\raggedright\ModernAnswerLabel\par\smallskip\centering
% Contract: planes y=0,1.4; grazing angle 30deg schematic; AC,AD normal to rays.
\begin{tikzpicture}[font=\small,>=Latex]
 \coordinate (A) at (0,1.4); \coordinate (B) at (0,0);
 \coordinate (C) at (-.606218,.35); \coordinate (D) at (.606218,.35);
 \foreach \y in {0,1.4}{\draw[TextGray] (-3.2,\y)--(3.4,\y);}
 \draw[vector] (-2.598076,2.9)--(A);
 \draw[vector] (A)--(2.598076,2.9);
 \draw[vector] (-2.598076,1.5)--(B);
 \draw[vector] (B)--(2.598076,1.5);
 \draw[guide] (A)--(C); \draw[guide] (A)--(D);
 \foreach \p in {A,B,C,D}{\fill (\p) circle (1.3pt);}
 \node[above=6pt] at (A) {$A$}; \node[below=7pt] at (B) {$B$};
 \node[left=7pt] at (C) {$C$}; \node[right=7pt] at (D) {$D$};
 \draw[<->] (3.8,0)--(3.8,1.4) node[midway,right=5pt] {$d$};
 \draw (1,1.4) arc[start angle=0,end angle=30,radius=1];
 \node at (1.3,1.72) {$\theta$};
 \node[above] at (-2.4,2.9) {入射方向};
 \node[above] at (2.35,2.9) {出射方向};
\end{tikzpicture}

\figcap{晶面间距为 $d$ 的布拉格衍射几何。相邻晶面反射波程差为 $2d\sin\theta$。}\label{ch32:fig:electron-diffraction}
\end{minipage}\end{center}
图~\ref{ch32:fig:electron-diffraction} 中 $\theta$ 是电子束与晶面的掠射角，而不是与晶面法线的夹角。两束平行波分别在上层的 $A$ 和下层的 $B$ 处散射。
从 $A$ 向下层入射线、出射线作垂线，垂足分别为 $C,D$。$AC$ 是入射等相位面的截线，$AD$ 用来比较出射相位；下层波多走的路程正是 $CB+BD$。
因为 $AB=d$，两个直角三角形分别给出 $CB=d\sin\theta$、$BD=d\sin\theta$。
两段相加得到波程差 $2d\sin\theta$；同相叠加要求此值等于整数倍波长。图中角度为几何示意，本题另由数据计算。


非相对论德布罗意波长
\[
 \lambda=\frac{1.227}{\sqrt{54.0}}\,\mathrm{nm}
 \approx\boxed{0.167\,\mathrm{nm}}.
\]
一级布拉格条件 $2d\sin\theta=\lambda$ 给出
\[
 \sin\theta=\frac{1.227/\sqrt{54.0}}{0.400}\approx0.4174,
 \qquad \boxed{\theta\approx24.7^\circ}.
\]
$\sin\theta<1$，一级衍射允许。加速电压远小于 $511\,\mathrm{kV}$，非相对论近似自洽；若电压接近该尺度须修正动量关系。
\end{xesolution}



\section{原子光谱与玻尔模型}

氢原子光谱满足里德伯关系（Rydberg formula）
\begin{equation}
 \frac1\lambda=R_H\left(\frac1{n_f^2}-\frac1{n_i^2}\right),
 \qquad n_i>n_f.
\end{equation}
玻尔模型（Bohr model）先近似把质子视为静止，记 $k_e=1/(4\pi\varepsilon_0)$。库仑力提供向心力，角动量取离散值：
\[
 \frac{k_ee^2}{r_n^2}=\frac{m_ev_n^2}{r_n},\qquad m_ev_nr_n=n\hbar.
\]
把 $v_n=n\hbar/(m_er_n)$ 代入第一式，约去公共因子，得
\[
 r_n=\frac{n^2\hbar^2}{m_ek_ee^2}=a_0n^2,\qquad v_n=\frac{k_ee^2}{n\hbar}.
\]
以无限远处电子与质子静止时的能量为零，势能 $V=-k_ee^2/r_n$，动能 $K=m_ev_n^2/2=k_ee^2/(2r_n)$，故总能量
\[
 E_n=K+V=-\frac{k_ee^2}{2r_n}=-\frac{m_ek_e^2e^4}{2\hbar^2n^2}.
\]
代入常数后得到
\begin{equation}
 r_n=a_0n^2,\qquad E_n=-\frac{13.6\,\mathrm{eV}}{n^2}.
\end{equation}
跃迁光子满足 $h\nu=E_i-E_f$。玻尔模型正确解释氢样离子能级的主要结构，却不能完整描述多电子原子、谱线强度、精细结构和自旋；
这些需要量子力学。

原子尺度上电磁相互作用的无量纲强度由精细结构常数
\begin{equation}
 \alpha=\frac{e^2}{4\pi\varepsilon_0\hbar c}\approx\frac1{137}
\end{equation}
刻画。玻尔半径可写成 $a_0=\hbar/(m_ec\alpha)$，轨道速率为 $v_n=\alpha c/n$。
约化康普顿波长 $\bar\lambda_C=\hbar/(m_ec)$ 与经典电子半径 $r_e=k_ee^2/(m_ec^2)$ 因而满足
\[
 \bar\lambda_C=\alpha a_0,\qquad r_e=\alpha\bar\lambda_C=\alpha^2a_0.
\]
这里 $r_e$ 是由能量尺度定义的长度，不表示电子是半径为 $r_e$ 的刚性球。精细结构能量相对主能级通常带有 $\alpha^2$ 量级的小因子。这种无量纲参数比单独比较带单位常数更能显示修正项的相对大小。

\begin{origexamplebox}{氢原子的电离与谱线}
处于 $n=3$ 能级的氢原子跃迁到 $n=2$。（1）求光子能量；（2）从 $n=3$ 状态电离至少需要多少能量？
\end{origexamplebox}
\begin{xesolution}
\ModernAnswerLabel\quad \[
 E_3=-\frac{13.6}{9}\,\mathrm{eV},\qquad E_2=-\frac{13.6}{4}\,\mathrm{eV}.
\]
发射光子能量
\[
 \boxed{E_\gamma=E_3-E_2=13.6\left(\frac14-\frac19\right)\mathrm{eV}
 \approx1.89\,\mathrm{eV}}.
\]
从 $n=3$ 到连续谱零能级所需
\[
 \boxed{E_{\rm ion}=0-E_3=\frac{13.6}{9}\,\mathrm{eV}\approx1.51\,\mathrm{eV}}.
\]
两问对应不同终态，不能把跃迁光子能量当成该状态的电离能。
\begin{center}\begin{minipage}{\linewidth}\centering
% Contract: E2=-3.4, E3=-1.511, continuum=0; arrow directions encode processes.
\begin{tikzpicture}[>=Latex,font=\small,y=.85cm]
 \draw[->] (-.55,-3.9)--(-.55,.85) node[above]{$E/\mathrm{eV}$};
 \foreach \y/\lab in {-3.4/{n=2},-1.5111/{n=3},0/{n\to\infty}}{
  \draw[thick] (0,\y)--(6,\y);
  \node[right=5pt] at (6,\y){$\lab$};
 }
 \node[left=4pt] at (-.55,-3.4){$-3.40$};
 \node[left=4pt] at (-.55,-1.5111){$-1.51$};
 \node[left=4pt] at (-.55,0){$0$};
 \draw[DeepBlue,very thick,->] (1.1,-1.5111)--(1.1,-3.4);
 \node[right=6pt,align=left] at (1.1,-2.45){发光\\$1.89\,\mathrm{eV}$};
 \draw[DeepBlue,very thick,->] (4,-1.5111)--(4,0);
 \node[left=6pt,align=right] at (4,-.7){吸能电离\\$1.51\,\mathrm{eV}$};
\end{tikzpicture}

\figcap{从同一 $n=3$ 能级出发的两个不同过程，能级高度按能量绘制。}\label{ch32:fig:hydrogen-two-processes}
\end{minipage}\end{center}
图~\ref{ch32:fig:hydrogen-two-processes} 左箭头向下，原子的能量减少，
光子获得正能量 $E_3-E_2$；右箭头向上，原子需要吸收正能量 $0-E_3$。
零能量线表示电子与质子在无限远处、相对动能为零的电离阈值，不是又一个有限半径的束缚轨道。
\end{xesolution}


\subsection{氢样离子、约化质量与能级极限}
只有一个电子、核电荷为 $+Ze$ 的氢样离子，可以把库仑常数因子中的 $e^2$ 换成 $Ze^2$。若核质量为 $M$，电子和核均绕质心运动，相对坐标的质量为约化质量
\[
 \mu=\frac{m_eM}{m_e+M}.
\]
相应的轨道半径（两粒子的间距）和能级为
\[
 r_n=\frac{m_e}{\mu}\frac{a_0n^2}{Z},\qquad
 E_n=-\frac{\mu}{m_e}\frac{13.6Z^2}{n^2}\,\mathrm{eV}.
\]
因此氢样氦离子 $\mathrm{He}^+$ 的基态电离能约为 $54.4\,\mathrm{eV}$；中性氦有两个电子，不能直接套这条单电子公式。对同一 $Z$ 的同位素，较重的核使 $\mu$ 更接近 $m_e$，同一跃迁的光子能量略增、波长略短。
这里 $13.6\,\mathrm{eV}$ 已取近似值，精密同位素位移还需一致使用常数精度及其他修正。玻尔近似的速度尺度为 $Z\alpha c/n$；当它不再远小于 $c$ 时，非相对论轨道模型不适用。

\subsection{氢光谱的级次}

\begin{origexamplebox}{巴耳末系最长波长与系限}
取里德伯常数 $R_H=1.097\times10^7\,\mathrm{m^{-1}}$。求氢原子巴耳末系（终态 $n_f=2$）的最长波长和短波极限。
\end{origexamplebox}
\begin{xesolution}
\begin{center}\begin{minipage}{\linewidth}\centering\raggedright\ModernAnswerLabel\par\smallskip\centering
% Contract: upper levels proportional in eV; explicitly broken axis below n=2.
\begin{tikzpicture}[font=\small,>=Latex]
 \draw[axis] (-2.6,-.35)--(-2.6,.7);
 \draw[axis] (-2.6,1.25)--(-2.6,6.15) node[above] {$E\;(\mathrm{eV})$};
 \draw (-2.72,.86)--(-2.48,1.00); \draw (-2.72,1.06)--(-2.48,1.20);
 \foreach \n/\y/\en in {1/0/-13.6,2/2.2/-3.40,3/4.089/-1.51,4/4.75/-0.85}{
  \draw[DeepBlue,thick] (-.1,\y)--(4.8,\y);
  \node[left=6pt] at (-.1,\y) {$n=\n$};
  \node[left=5pt] at (-2.6,\y) {$\en$};
 }
 \draw[TextGray,dashed] (-.1,5.6)--(4.8,5.6);
 \node[left=6pt] at (-.1,5.6) {$n=\infty$};
 \node[left=5pt] at (-2.6,5.6) {$0$};
 \draw[vector] (1.4,5.6)--(1.4,2.2);
 \draw[vector] (3.8,4.089)--(3.8,2.2);
 \node[below=8pt,align=center] at (1.4,2.2) {巴耳末系限};
 \node[below=8pt] at (3.8,2.2) {$3\to2$};
 \node at (2.4,.9) {纵轴断开};
\end{tikzpicture}

\figcap{氢原子能级与巴耳末系。纵轴在基态与第一激发态之间断开；上部各能级按能量比例绘制，越靠近电离极限，间隔越小。}\label{ch32:fig:hydrogen-levels}
\end{minipage}\end{center}


最长波长对应最小能级差，即 $n_i=3$：
\[
 \frac1{\lambda_{\max}}=R_H\left(\frac14-\frac19\right)
 =R_H\frac5{36},
\]
所以 $\boxed{\lambda_{\max}\approx656\,\mathrm{nm}}$。短波极限对应 $n_i\to\infty$：
\[
 \frac1{\lambda_{\min}}=\frac{R_H}{4},
 \qquad \boxed{\lambda_{\min}\approx364.6\,\mathrm{nm}}.
\]
图~\ref{ch32:fig:hydrogen-levels} 显示高能级向电离极限聚集，故跃迁光子能量趋于有限上限而不是无限增大。
\end{xesolution}



\section{原子核、结合能与核反应}

原子的典型尺度约为 $10^{-10}\,\mathrm m$；原子核半径通常在 $10^{-15}$--$10^{-14}\,\mathrm m$ 范围，
经验关系为
\begin{equation}
 R\simeq r_0A^{1/3},\qquad r_0\approx1.2\,\mathrm{fm}.
\end{equation}
核质量小于自由核子质量之和，质量亏损对应结合能
\begin{equation}
 B=\Delta mc^2.
\end{equation}
平均每核子结合能反映核的相对稳定性。轻核聚变和重核裂变都可能向更高平均结合能方向释放能量。


\subsection{质量表、反应Q值与两体反冲}
记核电荷数为 $Z$、中子数为 $N$、质量数为 $A=Z+N$。核反应前后先核对总电荷与总核子数，再计算质量差。若用裸核质量 $M_N$，结合能为
\[
 B=[Zm_p+(A-Z)m_n-M_N]c^2.
\]
若给的是中性原子质量 $M_{\rm at}$，并忽略电子结合能差，宜改用氢原子质量 $m_H$：
\[
 B\simeq[Zm_H+(A-Z)m_n-M_{\rm at}]c^2.
\]
两式中的质量对象必须一致；把原子质量和裸质子质量直接混在一起会多减去 $Z$ 个电子静质量。
一般核反应的 $Q=(\sum M_i-\sum M_f)c^2$，并满足 $K_f-K_i=Q$。$Q<0$ 时固定靶入射阈能通常大于 $|Q|$，因为末态还需携带总动量；精确计算用第31章的系统不变量。

当静止母核发生两体衰变，产物动量相反且大小相等。若两产物均非相对论，有 $K_1=p^2/(2m_1)$、$K_2=p^2/(2m_2)$，从而
\[
 K_1=Q\frac{m_2}{m_1+m_2},\qquad
 K_2=Q\frac{m_1}{m_1+m_2}.
\]
较轻产物带走较多动能。若一产物为光子，须用 $p_\gamma=E_\gamma/c$，不能把光子代入 $p^2/(2m)$。母体静质量 $M_i$、子体质量 $M_f$ 的两体辐射跃迁精确给出
\[
 E_\gamma=\frac{M_i^2-M_f^2}{2M_i}c^2<Q=(M_i-M_f)c^2,
\]
差值用于子体反冲。忽略反冲是建立在 $Q\ll M_ic^2$ 上的近似。

\subsection{结合能曲线与裂变能}

\begin{origexamplebox}{由平均结合能估算裂变能}
粗略取 $^{235}\mathrm U$ 的平均每核子结合能为 $7.6\,\mathrm{MeV}$，裂变产物平均为
$8.5\,\mathrm{MeV}$。忽略裂变中子与产物间的细小差别，估算一次裂变释放的能量，并换算为焦耳。
\end{origexamplebox}
\begin{xesolution}
\begin{center}\begin{minipage}{\linewidth}\centering\raggedright\ModernAnswerLabel\par\smallskip\centering
% Contract: schematic binding trend, arrows separate from approximate data.
\begin{tikzpicture}
\begin{axis}[width=.83\linewidth,height=5.8cm,axis lines=left,
 xlabel={$A$},ylabel={$B/A\;(\mathrm{MeV})$},xmin=0,xmax=250,ymin=0,ymax=11,
 xtick={56,120,235},ytick={0,4,8},every axis label/.style={fill=none}]
 \addplot[DeepBlue,very thick] coordinates{(1,0)(2,1.1)(4,7.1)(12,7.7)(20,8.0)(40,8.5)(56,8.8)(90,8.7)(120,8.5)(160,8.2)(200,7.9)(235,7.6)};
 \node[above=6pt] at (axis cs:56,8.8) {铁、镍附近};
 \draw[vector] (axis cs:220,5.7)--(axis cs:130,6.7);
 \node[below=6pt] at (axis cs:175,5.7) {重核裂变};
 \draw[vector] (axis cs:15,4.0)--(axis cs:50,6.1);
 \node[right=5pt] at (axis cs:42,4.5) {轻核聚变};
\end{axis}
\end{tikzpicture}

\figcap{平均每核子结合能的典型趋势。轻核聚变和重核裂变都可向更高结合能方向释放能量。}\label{ch32:fig:binding-curve}
\end{minipage}\end{center}
图~\ref{ch32:fig:binding-curve} 中箭头只指示反应趋向，不表示精确核素之间的能量差；折线为典型趋势，不能代替核质量表。本题用题给的反应前后平均结合能之差，乘参与反应的核子数估算释放能量。


总结合能增加约
\[
 Q\approx235(8.5-7.6)\,\mathrm{MeV}
 =\boxed{2.12\times10^2\,\mathrm{MeV}}.
\]
用 $1\,\mathrm{MeV}=1.602\times10^{-13}\,\mathrm J$，
\[
 \boxed{Q\approx3.39\times10^{-11}\,\mathrm J}.
\]
这是数量级估算，精确值应使用各核素静质量和实际产物分支。符号为正，因为末态总结合能更大、总静质量更小。
\end{xesolution}

\subsection{核反应的Q值}

\begin{origexamplebox}{氘氚聚变的Q值}
反应 $^2\mathrm H+{}^3\mathrm H\to{}^4\mathrm{He}+n$ 中，氘、氚、氦的中性原子质量以及中子的质量依次为
$2.014102\,\mathrm u$、$3.016049\,\mathrm u$、$4.002603\,\mathrm u$、$1.008665\,\mathrm u$。
取 $1\,\mathrm u\,c^2=931.5\,\mathrm{MeV}$，求Q值。
\end{origexamplebox}
\begin{xesolution}
\begin{center}\begin{minipage}{\linewidth}\centering\raggedright\ModernAnswerLabel\par\smallskip\centering
% Contract: reaction bookkeeping, not spatial particle trajectories.
\begin{tikzpicture}[font=\small,>=Latex]
 \node[draw,circle,minimum size=11mm] (d) at (-4,0) {$^2\mathrm H$};
 \node at (-2.95,0) {$+$};
 \node[draw,circle,minimum size=11mm] (t) at (-1.9,0) {$^3\mathrm H$};
 \draw[vector] (-.95,0)--(.95,0);
 \node[draw,circle,minimum size=13mm] (he) at (2,0) {$^4\mathrm{He}$};
 \node at (3.15,0) {$+$};
 \node[draw,circle,minimum size=8mm] (n) at (4.15,0) {$n$};
 \node at (-2.95,1.1) {反应物}; \node at (3.05,1.1) {产物};
 \node[align=center] at (0,-1.3) {$Q=(m_D+m_T-m_{\rm He}-m_n)c^2>0$\\总动能增加，静质量减少};
\end{tikzpicture}

\figcap{氘氚聚变的质量--能量收支。核子数和电荷数在反应前后分别守恒。}\label{ch32:fig:fusion-q}
\end{minipage}\end{center}
图~\ref{ch32:fig:fusion-q} 把反应物和产物分列两侧，下面计算先用质量差判断 $Q$ 值符号，再核验两个守恒数。


两侧电子数相同，电子静质量在相减时抵消；电子结合能仍有很小差异，
但在本题约 $17.6\,\mathrm{MeV}$ 的精度下可以忽略。因此用所给中性原子质量计算。质量差
\[
 \Delta m=(2.014102+3.016049-4.002603-1.008665)\,\mathrm u
 =0.018883\,\mathrm u.
\]
故
\[
 \boxed{Q=\Delta mc^2\approx0.018883\times931.5
 \approx17.6\,\mathrm{MeV}}.
\]
$Q>0$，静质量减少量转化为产物动能。核子数 $2+3=4+1$、电荷数 $1+1=2+0$ 均守恒，是写反应式时的第一层检查。
\end{xesolution}



\section{放射性衰变与粒子守恒}

放射性核数满足
\begin{equation}
 N(t)=N_0e^{-\lambda t},\qquad
 T_{1/2}=\frac{\ln2}{\lambda},\qquad
 \mathcal A(t)=\lambda N(t).
\end{equation}
活度记为 $\mathcal A$，单位为贝可勒尔（$\mathrm{Bq}=\mathrm{s^{-1}}$），以免与质量数 $A$ 混淆。指数规律来自 $\dd N/\dd t=-\lambda N$：分离变量并积分得 $\ln[N(t)/N_0]=-\lambda t$；令 $N/N_0=1/2$ 即得半衰期。半衰期是统计规律，不表示某个原子核在半衰期时刻必然衰变。


\begin{origexamplebox}{经过三个半衰期的剩余量}
某样品半衰期为 $T_{1/2}$。经过 $3T_{1/2}$ 后，未衰变核数和活度分别为初值的多少？
\end{origexamplebox}
\begin{xesolution}
\ModernAnswerLabel\quad 每经过一个半衰期剩余一半，故
\[
 \boxed{\frac N{N_0}=2^{-3}=\frac18}.
\]
同一核素的 $\lambda$ 不变，活度 $\mathcal A=\lambda N$，所以
\[
 \boxed{\frac{\mathcal A}{\mathcal A_0}=\frac18}.
\]
若衰变产物也有放射性，总测得活度还要加上子核贡献，不能只用单指数。
\end{xesolution}


\subsection{衰变方式与计数模型}
常见衰变方式可用守恒量区分：
\begin{align*}
 {}^A_ZX&\to{}^{A-4}_{Z-2}Y+{}^4_2\mathrm{He} &&(\alpha),\\
 {}^A_ZX&\to{}^A_{Z+1}Y+e^-+\bar\nu_e &&(\beta^-),\\
 {}^A_ZX&\to{}^A_{Z-1}Y+e^++\nu_e &&(\beta^+),\\
 {}^A_ZX^*&\to{}^A_ZX+\gamma &&(\gamma).
\end{align*}
$\beta$ 衰变中的电子或正电子是在相互作用中产生的，不是从核内预先存放的经典电子中取出。三体末态中能量在带电轻子、中微子与反冲核之间分配，所以电子能谱连续；这不表示每次衰变的总能量不守恒。$\gamma$ 跃迁不改变核的 $A,Z$。
若用中性原子质量并忽略电子结合能差，$\beta^-$ 衰变约有 $Q=[M_X-M_Y]c^2$，而 $\beta^+$ 衰变约有 $Q=[M_X-M_Y-2m_e]c^2$；后者额外需要两个电子静能的质量差门槛。不能把两式混用。

探测器的净计数率还与效率 $\varepsilon$ 有关：$R(t)=\varepsilon\mathcal A(t)$。若效率稳定、只计单一核素且背景已扣除，间隔 $\Delta t$ 的计数率比满足
\[
 \lambda=\frac1{\Delta t}\ln\frac{R(t)}{R(t+\Delta t)}.
\]
从 $t=0$ 到 $\Delta t$ 的期望净计数应积分：$C=\varepsilon N_0(1-e^{-\lambda\Delta t})$，只有 $\lambda\Delta t\ll1$ 时才可近似为初始计数率乘时间。任意区间 $[t,t+\Delta t]$ 应将上式的 $N_0$ 换为 $N(t)$。若有恒定背景，应先减去背景再取对数；若子核也能被计数，应先建立衰变链模型。

\subsection{两级放射性衰变链}

\begin{origexamplebox}{母核--子核衰变链的峰值}
母核1以衰变常数 $\lambda_1$ 变成放射性子核2，子核2再以 $\lambda_2=2\lambda_1$ 衰变为稳定核。
初始有 $N_{10}$ 个母核且无子核。求 $N_2(t)$、子核数达到最大值的时刻和最大值。
\end{origexamplebox}
\begin{xesolution}
\begin{center}\begin{minipage}{\linewidth}\centering\raggedright\ModernAnswerLabel\par\smallskip\centering
% Contract: lambda2=2lambda1, N2 peak (ln2,1/4), legend outside curves.
\begin{tikzpicture}
\begin{axis}[width=.8\linewidth,height=5.4cm,axis lines=left,
 xlabel={$\lambda_1t$},ylabel={$N/N_{10}$},xmin=0,xmax=5,ymin=0,ymax=1.08,
 xtick={.693147,2,4},xticklabels={$\ln2$,$2$,$4$},ytick={.25,.5,1},
 legend style={draw=none,fill=none,at={(.96,.94)},anchor=north east},
 every axis label/.style={fill=none}]
 \addplot[DeepBlue,very thick,domain=0:5,samples=200]{exp(-x)}; \addlegendentry{母核 $N_1$}
 \addplot[TextGray,very thick,dashed,domain=0:5,samples=200]{exp(-x)-exp(-2*x)}; \addlegendentry{子核 $N_2$}
 \draw[guide] (axis cs:.693147,0)--(axis cs:.693147,.25)--(axis cs:0,.25);
 \addplot[only marks,mark=*,DeepBlue,forget plot] coordinates {(.693147,.25)};
\end{axis}
\end{tikzpicture}

\figcap{$\lambda_2=2\lambda_1$ 时母核与子核数目的变化。子核先积累，达到峰值后再衰减。}\label{ch32:fig:decay-chain}
\end{minipage}\end{center}


母核数 $N_1=N_{10}e^{-\lambda_1t}$。子核满足
\[
 \dot N_2=\lambda_1N_1-\lambda_2N_2,
 \qquad N_2(0)=0.
\]
将方程写成 $\dot N_2+\lambda_2N_2=\lambda_1N_{10}e^{-\lambda_1t}$。两侧乘积分因子 $e^{\lambda_2t}$，左侧恰好成为乘积的导数：
\[
 \frac{\dd}{\dd t}(N_2e^{\lambda_2t})=\lambda_1N_{10}e^{(\lambda_2-\lambda_1)t}.
\]
从0积分到 $t$，利用 $N_2(0)=0$，当 $\lambda_2\ne\lambda_1$ 时有
\[
 N_2e^{\lambda_2t}=\lambda_1N_{10}\frac{e^{(\lambda_2-\lambda_1)t}-1}{\lambda_2-\lambda_1}.
\]
两侧乘 $e^{-\lambda_2t}$，得到一般的巴特曼形式
\[
 N_2=N_{10}\frac{\lambda_1}{\lambda_2-\lambda_1}
 \left(e^{-\lambda_1t}-e^{-\lambda_2t}\right).
\]
代入 $\lambda_2=2\lambda_1$：
\[
 \boxed{N_2=N_{10}\left(e^{-\lambda_1t}-e^{-2\lambda_1t}\right)}.
\]
求导得 $\dot N_2=\lambda_1N_{10}e^{-\lambda_1t}[-1+2e^{-\lambda_1t}]$。
前面的因子始终为正，方括号由正变负，故唯一零点就是最大值点。令 $e^{-\lambda_1t}=1/2$，所以
\[
 \boxed{t_{\max}=\frac{\ln2}{\lambda_1}},
 \qquad \boxed{N_{2,\max}=\frac{N_{10}}4}.
\]
$t=0$ 与 $t\to\infty$ 时 $N_2$ 都为零，峰值位于中间，和图~\ref{ch32:fig:decay-chain} 一致。
\end{xesolution}

\subsection{衰变链的相等常数极限}
一般两级模型中令 $\lambda_2\to\lambda_1=\lambda$，巴特曼式表面上为 $0/0$，应先取极限而非宣布发散。用导数或直接积分得
\[
 N_2(t)=\lambda N_{10}t e^{-\lambda t},\qquad
 t_{\max}=\frac1\lambda,\qquad N_{2,\max}=\frac{N_{10}}e.
\]
无论两常数是否相等，子核峰值时都满足生产率等于损失率：$\lambda_1N_1=\lambda_2N_2$。相等常数极限、初始导数 $\dot N_2(0)=\lambda_1N_{10}$、长时间趋零，是检查衰变链公式的三个独立办法。若母核只有比例 $b$ 的分支通向所计子核，源项应改为 $b\lambda_1N_1$；若初始已有子核，还需加 $N_{20}e^{-\lambda_2t}$。

\subsection{中微子与粒子物理概览}

连续的 $\beta$ 衰变能谱促使泡利（Wolfgang Pauli）提出中微子，以保持逐次事件中的能量、动量和角动量守恒。中微子只参与弱相互作用和引力，
探测率极低。2002年诺贝尔物理学奖的一半授予戴维斯和小柴昌俊，表彰宇宙中微子探测方面的开创性工作；另一半授予贾科尼，表彰宇宙X射线源研究。2015年奖项表彰中微子振荡的发现，
这一现象表明中微子具有非零质量。两次奖项对应不同成果，不能合并成一次“发现中微子质量”的奖项。

标准模型把物质基本粒子分为夸克和轻子，以规范玻色子传递强、电磁和弱相互作用。质子、中子由夸克和胶子组成；
夸克具有分数电荷但在通常条件下不能孤立观测。标准模型不包含完整的量子引力理论，也不能单独解释暗物质等所有现象。




\hypersetup{bookmarksdepth=2}
\endgroup
